\documentclass[11pt]{article}
% Shared print layout. The web's provisional design and MathJax are independent.
\usepackage[utf8]{inputenc}
\usepackage[T1]{fontenc}
\usepackage[spanish,es-nodecimaldot]{babel}
\usepackage{lmodern}
\usepackage{amsmath,amssymb}
\usepackage{graphicx}
\usepackage{geometry}
\geometry{a4paper,left=20mm,right=20mm,top=25mm,bottom=22mm,headheight=16pt,headsep=8mm}
\usepackage{fancyhdr,lastpage}
\usepackage{enumitem}
\usepackage{xcolor}
\usepackage{pgfplots}
\pgfplotsset{compat=1.18}
\usepgfplotslibrary{fillbetween}
\usepackage[hidelinks,pdfencoding=auto]{hyperref}
\hypersetup{pdflang=es-ES,pdfauthor={},pdfsubject={Matemáticas Aplicadas a las Ciencias Sociales II}}
\pagestyle{fancy}
\fancyhf{}
\fancyhead[L]{\small Matemáticas Aplicadas a las Ciencias Sociales II}
\fancyhead[R]{\small 2.\textsuperscript{o} Bachillerato}
\fancyfoot[L]{\footnotesize Repaso de Análisis | Edición revisada 1.0}
\fancyfoot[R]{\footnotesize Página \thepage\ de \pageref{LastPage}}
\renewcommand{\headrulewidth}{0.4pt}
\setlength{\parindent}{0pt}
\setlength{\parskip}{0.5em}
\setlist{itemsep=0.3em,topsep=0.3em}
\newcommand{\apartado}[1]{\subsection*{#1}}
\pgfplotsset{
  revision/.style={
    width=0.94\linewidth,height=7.5cm,
    axis lines=middle,grid=major,grid style={gray!18},
    tick label style={font=\small},label style={font=\small},
    xlabel={$x$},ylabel={$y$},samples=160,
    legend style={at={(0.5,-0.16)},anchor=north,draw=none,font=\small},
    clip=true
  }
}

\hypersetup{pdftitle={Repaso de Análisis - Exámenes 1 y 2 - Examen revisado}}
\begin{document}
{\Large\bfseries Repaso de Análisis: dos exámenes}

Se conservan ambos exámenes de 10 puntos por separado. El original no fija duración ni autorización de calculadora; no se importan condiciones de otros exámenes. En derivadas y límites se pide elegir cuatro y solo hay cuatro opciones: se conservan todas.

Apellidos y nombre: \hrulefill

Fecha: \hrulefill


\section*{Examen 1: 10 puntos}
\subsection*{Ejercicio 1 \hfill 3,2 puntos} Se considera la función:
\[
f(x) =
\begin{cases} 
x^2 + ax + 2, & \text{si } x \leq 0, \\
\dfrac{x + b}{x - 1}, & \text{si } x > 0,
\end{cases}
\]
con \( a \) y \( b \) reales.

\begin{enumerate}[label=(\alph*)]
\item (2 puntos)  \textquestiondown{}Para qué valores de \(a\) y \(b\) la función es continua y derivable en \(x = 0\)?



\item (0,6 puntos)  Calcula la recta tangente a \( f(x) \) en \( x = 2 \).



\item (0,6 puntos)  Calcula la recta normal a \( f(x) \) en \( x = 2 \).



\end{enumerate}

\subsection*{Ejercicio 2 \hfill 2,4 puntos} Calcula las derivadas de cuatro de las siguientes funciones:


\begin{enumerate}[label=(\alph*)]
\item \( f(x) = \left( \frac{3x^2 - 5}{2x + 3} \right)^3 \)



\item \( f(x) = \sin^2{x} \)



\item \( f(x) = \frac{x^3 - 4}{(x^2 + 4)^2} \)



\item \( f(x) = e^{3x^2 - x + 2} \)


\end{enumerate}


\clearpage
\subsection*{Ejercicio 3 \hfill 2,4 puntos} Calcula cuatro de los siguientes límites:


\begin{enumerate}[label=(\alph*)]
\item \( \displaystyle \lim_{x \to 0} \frac{1 - \sqrt{1 - x^2}}{x^2} \)



\item \( \displaystyle \lim_{x \to +\infty} (3^x - x^5) \)



\item \( \displaystyle \lim_{x \to +\infty} \frac{5x + 7}{\sqrt{2x^2 + 3}} \)



\item \( \displaystyle\lim_{x\to+\infty}\left(\frac{x^2-1}{x^2}\right)^x\) (tendencia explicitada según la solución original; faltaba en el enunciado).


\end{enumerate}


\subsection*{Ejercicio 4 \hfill 2 puntos} La distancia entre un cohete y la zona de aterrizaje está dada por:
\[
D(t) = \frac{100}{t^2 + 1}
\]

\begin{enumerate}[label=(\alph*)]
\item (1 punto)  \textquestiondown{}A qué distancia se encuentra al inicio y tras 2 horas?



\item (1 punto)  \textquestiondown{}A qué valor tiende la distancia cuando \( t \to \infty \)?


\end{enumerate}






\clearpage
\section*{Examen 2: 10 puntos}
\subsection*{Ejercicio 1 \hfill 2 puntos} Se considera la función:
\[
f(x) = x^3 - 3x^2 + 2x.
\]
\begin{enumerate}[label=(\alph*)]
\item (1 punto) Halle los puntos de corte con los ejes, los intervalos de crecimiento y decrecimiento, los extremos relativos de \(f\) y su curvatura.



\item (1 punto) Represente gráficamente la función \(f\).


\end{enumerate}

\subsection*{Ejercicio 2 \hfill 3 puntos} Se considera la función:
\[
f(x) =
\begin{cases} 
\frac{x^2}{3}, & 0 \leq x \leq 2 \\
\frac{4}{x+1}, & x > 2
\end{cases}
\]

\begin{enumerate}[label=(\alph*)]
\item (1.5 puntos) Estudie la continuidad y derivabilidad de la función \(f\).



\item (1.5 puntos) Determine los intervalos de crecimiento y decrecimiento, el máximo de la función y represente gráficamente la función \(f\).


\end{enumerate}

\clearpage
\subsection*{Ejercicio 3 \hfill 2 puntos} Determine las asíntotas de la función:
\[
f(x) = \frac{1}{(x-2)(x-1)}
\]



\subsection*{Ejercicio 4 \hfill 3 puntos} La temperatura en el interior de un equipo de refrigeración durante un día que sufrió un corte de energía viene dada por la función:

\[
f(t) =
\begin{cases} 
-9, & 0 \leq t \leq 1 \\
-t^2 + 12t - 20, & 1 < t < 11 \\
-9, & 11 \leq t \leq 24
\end{cases}
\]

\begin{enumerate}[label=(\alph*)]
\item (1 punto) Estudie la continuidad de \(f\).



\item (1 punto) Represente gráficamente la función \(f\).



\item (0.5 puntos) \textquestiondown{}A qué hora se produjo el corte de energía y cuánto duró?



\item (0.5 puntos) \textquestiondown{}Se estropearon los sueros o las vacunas?


\end{enumerate}



\end{document}
